%% USPSC-Resumo.tex
\setlength{\absparsep}{18pt} % ajusta o espaçamento dos parágrafos do resumo
\begin{resumo}
	\begin{flushleft}
			\setlength{\absparsep}{0pt} % ajusta o espaçamento da referência
			\SingleSpacing
			\imprimirautorabr~~\textbf{\imprimirtituloresumo}.	\imprimirdata. \pageref{LastPage} p.
			\imprimirtipotrabalho~-~\imprimirinstituicao, \imprimirlocal, \imprimirdata.
 	\end{flushleft}
\OnehalfSpacing
Este trabalho aborda a otimização de baixo nível de kernels de inferência aplicados a modelos de visão computacional, com foco na arquitetura Vision Transformer (ViT), motivada pelas restrições de latência e energia da percepção embarcada em Veículos Aéreos Não Tripulados (VANTs). O objetivo principal é reduzir o gargalo de transferência de dados entre a memória global e os Streaming Multiprocessors (SM) da GPU, fenômeno associado à ``Memory Wall'', particularmente severo no hardware de borda que equipa esses veículos. Para isso, desenvolve-se um kernel de atenção fundido em OpenAI Triton, explorando tiling, reuso de dados em registradores e SRAM, e cálculo de softmax online para evitar a materialização completa da matriz intermediária de atenção. A validação compara a implementação customizada com baselines em PyTorch por meio de latência, throughput, erro numérico, tráfego de memória e análise pelo Modelo Roofline. Os resultados indicam que a fusão de operadores reduz substancialmente o tráfego em memória global e acelera a execução em relação à implementação explícita, embora permaneça abaixo do backend otimizado do PyTorch. A partir das intensidades aritméticas medidas, discutem-se as implicações para plataformas embarcadas típicas de VANTs, nas quais a menor banda de memória torna a fusão de operadores proporcionalmente mais vantajosa. O trabalho contribui com uma implementação experimental e um protocolo de análise para investigar gargalos de desempenho em kernels de atenção aplicados a Vision Transformers.

\textbf{Palavras-chave}: Visão Computacional. Vision Transformers. Otimização de Kernels. OpenAI Triton. Engenharia de Performance. Computação Embarcada. VANTs.
\end{resumo}
